% Unidad U008. Clasificación orbital y elevación; redacción autónoma.

\chapter{Orbital classification, initial words and 1st boundary}
\label{chap:orbitas-elevacion-r36}

The ternary reduction of the emitter realizes 243 words within the ambient
fiber \(\mathbb F_3^6\). This unit does not choose three words by
decimal comparison. It first constructs the symmetry action of the catalogue,
classifies its orbits by integer invariants, and only afterwards orients
the three structural representatives.

\section{Dihedral action on \(6\) coordinates}

For \(u=(u_1,u_2,u_3,u_4,u_5,u_6)\in\mathbb F_3^6\), define
\begin{align*}
 r(u)&=(u_3,u_1,u_2,u_5,u_6,u_4),\\
 s(u)&=(u_4,u_5,u_6,u_1,u_2,u_3).
\end{align*}
We verify
\[
 r^3=s^2=1,
 \qquad srs=r^{-1}.
\]
Therefore, \(r,s\) generate an action of \(D_3\).

The same permutation acts on the six blocks of each emission
\(U=(U_1,\ldots,U_6)\). The projection \(\Psi\) satisfies
\[
 \Psi(gU)=g\Psi(U),
 \qquad g\in D_3.
\]

\begin{theorem}[Orbital decomposition of the catalog]
\label{thm:catalogo-descomposicion-d3}
The 243 realized words decompose into 43 orbits:
\[
 38\text{ orbits of cardinality }6,
 \qquad
 5\text{ orbits of cardinality }3.
\]
In particular,
\[
 38\cdot6+5\cdot3=243.
\]
\end{theorem}

\begin{proof}
The set is applied to each word.
\[
 \{u,ru,r^2u,su,sru,sr^2u\}
\]
and repetitions are removed. Exhaustive enumeration of the catalog produces
the indicated histogram. The equivariance of \(\Psi\) ensures that no
orbit leaves the realized image.
\end{proof}

\section{Fiber Invariants}

For a word realized \(w\), let
\[
 \mathcal F_w:=\{U\in\operatorname{im}T_{\min}:\Psi(U)=w\},
\]
\[
 n(w):=|\mathcal F_w|,
 \qquad
 M(w):=\sum_{U\in\mathcal F_w}\operatorname{mult}(U),
\]
and
\[
 c(w)=\bigl(\#\{j:w_j=0\},\#\{j:w_j=1\},\#\{j:w_j=2\}\bigr).
\]
The enriched catalogue further preserves, for each emission, a diagonal
coordinate \(\operatorname{diag}_c(U)\in\{0,\ldots,8\}\), a cardinal
orientation \(h_+(U)\), and the multiplicative signature \(E_{\rm sig}(U)\). All
of them are computed from the seed route; they are not read from a target
constant. Their componentwise algorithm will be printed alongside the complete
catalogue in the certification appendix and forms part of the input state
of the following predicates.

\begin{definition}[Closure orbit]
\label{def:orbita-clausura}
An orbit \(\mathcal O\) is of closure if it has cardinality six and, for
each \(w\in\mathcal O\),
\[
 n(w)=5,
 \qquad
 M(w)=1008,
\]
\[
 \{\operatorname{mult}(U):U\in\mathcal F_w\}
 =\{432,144,144,144,144\},
\]
and the common tritic census is \(c(w)=(2,3,1)\).
\end{definition}

\begin{definition}[Radical Minimal Orbits]
\label{def:orbitas-minimas-radicales}
An orbit \(\mathcal O\) is minimally radical if it has cardinality six, each
word possesses a unique emission of multiplicity \(144\), and
\[
 \{\operatorname{diag}_c(U):\Psi(U)\in\mathcal O\}=\{0,3,6\}.
\]
Among them, the orbit with census \((2,3,1)\) is called the propagation
orbit, and the orbit with census \((2,2,2)\) is called the self-scaling orbit.
\end{definition}

\begin{theorem}[Uniqueness of the Three Structural Orbits]
\label{thm:unicidad-tres-orbitas-estructurales}
In the catalogue of 468 emissions there exists a unique closure orbit, a
unique minimally radical orbit of propagation, and a unique minimally radical
orbit of self-scale.
\end{theorem}

\begin{proof}
The 43 orbits of \cref{thm:catalogo-descomposicion-d3} are enumerated. For each one, \(n(w)\), \(M(w)\), the multiset of multiplicities, the census \(c(w)\), and the diagonal support are computed. The predicates of \cref{def:orbita-clausura,def:orbitas-minimas-radicales} leave exactly one orbit in each class. The proof is exhaustive over a finite set and does not use decimal digits of the analytic characters.
\end{proof}

The three orbits are
\begin{align*}
\mathcal O_{\rm cl}
 &=\{001112,010211,100121,112001,121100,211010\},\\
\mathcal O_{\rm prop}
 &=\{011120,012110,101201,110012,120011,201101\},\\
\mathcal O_{\rm aut}
 &=\{002112,020211,112002,121200,200121,211020\}.
\end{align*}

\section{Internal Orientation of Each Orbit}

We fix the internal caliber.
\[
 \operatorname{diag}_c(U)=6,
 \qquad h_+(U)=ES.
\]
A word in an orbit is admissible in that gauge if some emission in
its fiber satisfies both conditions.

\begin{proposition}[Unique oriented representative]
\label{prop:representante-orientado-unico}
Each of the three structural orbits contains a unique word
admissible in the gauge \((6,ES)\). They are
\[
 \boxed{
 w_6^\pi=010211,
 \qquad
 w_6^e=201101,
 \qquad
 w_6^\varphi=121200.}
\]
\end{proposition}

\begin{proof}
The enumeration of the catalogue fibers counts the words containing
an emission with \(\operatorname{diag}_c=6\) and \(h_+=ES\). In each of
the three orbits, the number of coincidences is one. The superscripts
\(\pi,e,\varphi\) henceforth name the structural functions of
closure, propagation, and self-scaling; they did not participate in the selection.
\end{proof}

\section{Closure microfiber}

The fiber of the representative \(w_6^\pi\) consists of five emissions:
\[
\begin{array}{c|c|c|c|c}
U_6&E_{\rm sig}&\operatorname{mult}&\operatorname{diag}_c&\text{stratum}\\\hline
501|614|498|169|272|272&555555&432&4&\perp\\
810|923|870|169|272|272&339555&144&6&++\\
810|983|810|169|272|272&393555&144&6&++\\
870|923|810|169|272|272&933555&144&6&++\\
870|923|810|418|674|521&933717&144&5&--
\end{array}
\]
Thus,
\[
 432+4\cdot144=1008,
 \qquad
 5=1_\perp+3_{++}+1_{--}.
\]
The unipolar emissions oriented of the other two channels are
\[
 U_e=850|963|890|149|252|212,
\]
\[
 U_\varphi=830|943|830|109|252|252.
\]
The mass \(1008\) of this microfiber is not the period of the calendar
\(C_{108}\); the two numbers belong to different domains.

\section{Ternary lifting matrices}

On row vectors of \(\mathbb F_3^6\), the two unique transports
reconstructed from the oriented register produced by the TPK transition
are
\[
L_0=
\begin{pmatrix}
2&2&2&1&2&1\\
2&1&2&2&1&1\\
1&0&0&1&0&2\\
0&0&1&1&1&0\\
2&2&2&0&2&1\\
2&1&2&1&0&0
\end{pmatrix},
\qquad
L_1=
\begin{pmatrix}
2&1&2&0&2&2\\
1&2&1&1&2&0\\
1&2&1&1&1&2\\
0&1&0&0&2&1\\
2&0&2&1&0&1\\
0&2&2&2&1&1
\end{pmatrix}.
\]
The simultaneous compatibility of the three biographies uniquely selects
the lifting calendar \((L_0,L_0,L_1,L_1)\) among the sixteen
binary calendars of length four.

\begin{definition}[Block recurrence and concatenation]
\label{def:recurrencia-bloques-l0-l1}
For each channel \(\chi\in\{\pi,e,\varphi\}\), let
\(b_0^\chi=w_6^\chi\) and
\[
 b_{k+1}^\chi=b_k^\chi L_{\epsilon_k}\pmod3,
 \qquad
 (\epsilon_0,\epsilon_1,\epsilon_2,\epsilon_3)=(0,0,1,1).
\]
The accumulated word is
\[
 w_{30}^\chi
 =b_0^\chi\Vert b_1^\chi\Vert b_2^\chi
  \Vert b_3^\chi\Vert b_4^\chi.
\]
\end{definition}

This formulation distinguishes the current block \(b_k\in\mathbb F_3^6\) from
the concatenated word \(w_{6(k+1)}\). A word of length \(6k\) is not
multiplied by a matrix \(6\times6\).

\begin{theorem}[Ternary biographies of thirty positions]
\label{thm:palabras-w30}
The recurrence of the \cref{def:recurrencia-bloques-l0-l1} produce
\begin{align*}
w_{30}^{\pi}
 &=010211|012222|010211|002111|110221,\\
w_{30}^{e}
 &=201101|121221|102011|012222|102011,\\
w_{30}^{\varphi}
 &=121200|112202|121020|010210|010200.
\end{align*}
\end{theorem}

\begin{proof}
The five row vectors are successively multiplied by the indicated
matrices, reducing each coordinate modulo three after each product.
The concatenation of the five computed blocks gives the three expressions.
Each equality can be verified component by component by means of 24
scalar products per channel.
\end{proof}

The environmental indices of the five blocks are
\[
\begin{array}{c|rrrrr}
&b_0&b_1&b_2&b_3&b_4\\\hline
\pi&103&161&103&67&349\\
e&523&457&301&161&301\\
\varphi&450&398&438&102&99
\end{array},
\]
where each word of six trits is interpreted as an integer in base three.

\section{1st common boundary}

The boundary is obtained through a finite relation that preserves axis, aggregate,
saturation, column occupancy, and dual neutrality. The incidence matrix
used in this relation is
\begin{equation}
 A_W=
 \begin{pmatrix}
 0&1&1&1&1&1\\
 1&0&1&2&2&1\\
 1&1&0&1&2&2\\
 1&2&1&0&1&2\\
 1&2&2&1&0&1\\
 1&1&2&2&1&0
 \end{pmatrix}\in M_6(\mathbb F_3).
 \label{eq:u008-aw-frontera}
\end{equation}
If \(u_4^\chi\) is the fifth block of \(w_{30}^\chi\), the axis and the
aggregate prior to separating the two lateral channels are
\begin{equation}
 u_5^\varphi=u_4^e=102011,
 \qquad
 u_5^\pi+u_5^e=210112,
 \qquad
 u_4^\varphi A_WL_1^3=111101.
 \label{eq:u008-eje-agregado-r36}
\end{equation}
From them, they calculate
\begin{equation}
 q=012120,\qquad a=111101,\qquad
 qA_W=012000,\qquad aA_W=120210.
 \label{eq:u008-q-a-r36}
\end{equation}
Depth six imposes the saturation \(c=111111\); the minimal compatible
occupancy is \(\operatorname{colw}=223232\); and dual neutrality requires
\begin{equation}
 (BA_W)\mathbf1=000.
 \label{eq:u008-neutralidad-dual-r36}
\end{equation}

\begin{proposition}[Enumeration of the first boundary]
\label{prop:u008-enumeracion-r36}
The conditions \eqref{eq:u008-eje-agregado-r36}--
\eqref{eq:u008-neutralidad-dual-r36} leave exactly two matrices:
\begin{equation}
 B_-=
 \begin{pmatrix}021222\\222220\\102011\end{pmatrix},
 \qquad
 B_+=
 \begin{pmatrix}222220\\021222\\102011\end{pmatrix}.
 \label{eq:u008-dos-fronteras}
\end{equation}
The mirror involution swaps its first two rows. Its internal cards are, respectively,
\begin{equation}
 B_-\longmapsto(789,048,894),
 \qquad
 B_+\longmapsto(793,045,894).
 \label{eq:u008-cartas-frontera}
\end{equation}
The orientation APP of the cylinder semiabierto selecciona \(B_+\).
\end{proposition}

\begin{proof}
The axis fixes the third row and the aggregate fixes the sum of the first two.
Saturation and occupation eliminate the distributions that do not fill the
six columns. The enumeration of the solutions of the six linear equations
of \eqref{eq:u008-neutralidad-dual-r36} leaves the two matrices of
\eqref{eq:u008-dos-fronteras}; the semi-open chart and the orientation
distinguish the second without consulting a target decimal expansion.
\end{proof}

We therefore define,
\begin{equation}
 R_{36}:=B_+=
 \begin{pmatrix}
 222220\\
 021222\\
 102011
 \end{pmatrix},
 \qquad
 \operatorname{ind}_3(R_{36})=(726,215,301).
 \label{eq:u008-r36-construida}
\end{equation}
The object \(R_{36}\) contains three ordered blocks of six trits. It is not the
predictive quotient \(C_{36}^{\rm pred}\) of
\cref{thm:tpk-cociente-predictivo-36} or the excess of 36 classes of the seed
gauge.

\section{Residue--quotient lifting}

Let \(A\in M_6(\mathbb Z)\) be the integer lift of \(L_0\):
\[
 A=
\begin{pmatrix}
2&2&2&1&2&1\\
2&1&2&2&1&1\\
1&0&0&1&0&2\\
0&0&1&1&1&0\\
2&2&2&0&2&1\\
2&1&2&1&0&0
\end{pmatrix},
\qquad \det A=1.
\]
For \(x_k\in\mathbb Z^6\), define
\[
 y_k=x_kA,
 \qquad
 u_k=y_k\bmod3\in\{0,1,2\}^6,
 \qquad
 x_{k+1}=\frac{y_k-u_k}{3}.
\]
Then
\begin{equation}
 x_kA=u_k+3x_{k+1}.
\label{eq:bockstein-hensel-local}
\end{equation}

\begin{theorem}[Reversibility of the lift]
\label{thm:reversibilidad-bockstein-hensel}
The pair \((u_k,x_{k+1})\) is uniquely determined by \(x_k\), and
\[
 x_k=(u_k+3x_{k+1})A^{-1}.
\]
\end{theorem}

\begin{proof}
Euclidean division by three in each coordinate of \(x_kA\) determines
residue and quotient. Since \(\det A=1\), the inverse \(A^{-1}\) has
integer coefficients. Multiplying \cref{eq:bockstein-hensel-local} by it
recovers \(x_k\).
\end{proof}

In channel \(e\), the visible block \(102011\) reappears twice, but
the prestates modulo 27 are
\[
 (25,11,7,17,8,16)
 \neq
 (7,8,4,23,26,7),
\]
and the subsequent quotients are
\[
 (6,13,9,2,13,1)
 \neq
 (15,0,3,19,23,25).
\]
It is thus explicitly verified that the visible-phase return of a word is not the return of the enriched state.

\section{Obstructions to orbital selection and to the lifting \(w_6\to w_{30}\to R_{36}\): uniqueness, gauge, and Hensel reversibility}

\begin{criterion}[Orbital selection refutation criteria]
The unit is refuted if:
\begin{enumerate}
 \item the action \(D_3\) leaves the catalogue or changes multiplicities;
 \item any one of the three structural predicates leaves more than one orbit;
 \item the gauge \((6,ES)\) selects two words in an orbit;
 \item a target digit is used to define the predicates;
 \item the accumulated word \(w_{6k}\) is multiplied by \(L_0\) or
       \(L_1\) instead of the block \(b_k\);
 \item \(\det A\neq1\) or fails the inverse of the lift;
 \item two equal visible residues are declared to be equal complete states.
\end{enumerate}
\end{criterion}

The first joint boundary has already been constructed. The next unit
will classify the nine phase components, prove the nonadic return
without resetting, and construct the oriented mirror with its local resolution
\(3\to1\).
